% Source: https://artofproblemsolving.com/community/c3h1613930_geo_thing
% Stable identity: aops-topic-1613930-q01
Use coordinates $A=(0,0)$, $B=(2,0)$, $C=(2,2)$, $D=(0,2)$, so $K=(1,0)$, $L=(2,1)$, $M=(1,2)$, $N=(0,1)$.

\textbf{The point $P$.} Line $LA$ is $y=\dfrac {x}{2}$, and line $KC$ is $y=2(x-1)$. They meet where $\dfrac {x}{2}=2x-2$, i.e. $x=\dfrac{4}{3}$, so $P=\left(\dfrac{4}{3},\dfrac{2}{3}\right)$. By the symmetry of the square (a quarter turn about the centre $O=(1,1)$ maps the figure to itself and carries $K\to L\to M\to N$), $Q,R,T$ are the images of $P$, so the octagon is made of four congruent pieces $OKPL$, $OLQM$, $OMRN$, $ONTK$.

\begin{asy}
usepackage("amsmath");
size(5.5cm);
defaultpen(fontsize(10pt));
pair A=(0,0),B=(2,0),C=(2,2),D=(0,2),O=(1,1);
pair K=(1,0),L=(2,1),M=(1,2),N=(0,1);
pair P=(4/3,2/3),Q=(4/3,4/3),R=(2/3,4/3),T=(2/3,2/3);
fill(O--K--P--L--cycle,rgb(1,0.88,0.85));
fill(K--P--L--Q--M--R--N--T--cycle,rgb(0.88,0.92,1));
fill(O--K--P--L--cycle,rgb(1,0.88,0.85));
draw(A--B--C--D--cycle,linewidth(0.9));
draw(K--P--L--Q--M--R--N--T--cycle,blue+linewidth(0.9));
draw(O--K^^O--L^^O--M^^O--N,dashed);draw(O--P,dashed);
dot(O);label("$O$",O,NW);
label("$K$",K,S);label("$L$",L,E);label("$M$",M,N);label("$N$",N,W);label("$P$",P,(0.7,-0.9));
label("$A$",A,SW);label("$B$",B,SE);label("$C$",C,NE);label("$D$",D,NW);
label("\mbox{}",(min(currentpicture,user=true)+max(currentpicture,user=true))/2,black); // keeps the text after the figure black when LaTeX typesets the labels inline
\end{asy}

\textbf{One piece.} Quadrilateral $OKPL$ splits along $OP$ into triangles $OKP$ and $OPL$. With $O$ as origin, $\overrightarrow{OK}=(0,-1)$, $\overrightarrow{OP}=\left(\dfrac{1}{3},-\dfrac{1}{3}\right)$, $\overrightarrow{OL}=(1,0)$, so
\[
[OKP]=\dfrac{1}{2}\left|0\cdot\left(-\dfrac{1}{3}\right)-(-1)\cdot\dfrac{1}{3}\right|=\dfrac{1}{6},\qquad
[OPL]=\dfrac{1}{2}\left|\dfrac{1}{3}\cdot0-\left(-\dfrac{1}{3}\right)\cdot1\right|=\dfrac{1}{6} .
\]
So $[OKPL]=\dfrac{1}{3}$, and the octagon has area
\[
4\cdot\dfrac{1}{3}=\boxed{\dfrac{4}{3}}.
\]
